\documentclass[../../../../main.tex]{subfiles}
\begin{document}

$r$ は $0<r<1$ をみたす実数，$n$ は2以上の整数とする．
平面上に与えられた1つの円を，
次の条件 (i)，(ii) をみたす2つの円で置き換える操作 (P) を考える．

\begin{enumerate}
  \item 新しい2つの円の半径の比は $r:1-r$ で，半径の和はもとの円の半径に等しい．
  \item 新しい2つの円は互いに外接し，もとの円に内接する．
\end{enumerate}

以下のようにして，平面上に $2^n$ 個の円を作る．
\begin{itemize}
  \item 最初に，平面上に半径1の円を描く．
  \item 次に，この円に対して操作 (P) を行い，2つの円を得る(これを1回目の操作という)．
  \item $k$ 回目の操作で得られた $2^k$ 個の円のそれぞれについて，
  操作 (P) を行い，$2^{k+1}$ 個の円を得る $(1 \leqq k \leqq n-1)$．
\end{itemize}

\begin{figure}[htbp]
  \centering
  \begin{tikzpicture}[scale=0.8, >=stealth]
    \draw[thick] (0, 0) circle (1.5);

    \draw[->, thick] (2, 0) -- (3.5, 0) node[midway, above, align=center] {1回目\\の操作};

    \draw[thick] (5.5, 0.6) circle (0.9);
    \draw[thick] (5.5, -0.9) circle (0.6);

    \draw[->, thick] (7.2, 0) -- (8.7, 0) node[midway, above, align=center] {2回目\\の操作};

    \draw[thick] (10.5, 1.0) circle (0.5);
    \draw[thick] (10.5, 0.2) circle (0.3);
    \draw[thick] (10.5, -0.4) circle (0.3);
    \draw[thick] (10.5, -1.0) circle (0.3);

    \draw[->, thick] (11.8, 0) -- (12.8, 0);
    \node at (13.5, 0) {$\dots$};
  \end{tikzpicture}
\end{figure}

\begin{enumerate}
  \item $n$ 回目の操作で得られる $2^n$ 個の円の周の長さの和を求めよ．
  \item 2回目の操作で得られる4つの円の面積の和を求めよ．
  \item $n$ 回目の操作で得られる $2^n$ 個の円の面積の和を求めよ．
\end{enumerate}

\end{document}
